\documentclass[11pt,a4paper]{article}
\usepackage[T1]{fontenc}
\usepackage{lmodern}
\usepackage[margin=26mm]{geometry}
\usepackage{amsmath,amssymb,amsthm}
\usepackage{microtype}
\usepackage{enumitem}
\usepackage[hidelinks]{hyperref}
\setlength{\parskip}{3pt}
\setlength{\emergencystretch}{3em}
\newtheorem{theorem}{Theorem}[section]
\newtheorem{lemma}[theorem]{Lemma}
\newtheorem{corollary}[theorem]{Corollary}
\theoremstyle{remark}
\newtheorem{remark}[theorem]{Remark}
\newcommand{\Z}{\mathbb Z}
\newcommand{\N}{\mathbb N}
\newcommand{\Q}{\mathbb Q}
\newcommand{\F}{\mathbb F}
\newcommand{\ps}[1]{#1[[x]]}
\hypersetup{pdftitle={A396093: A complete parity classification},pdfauthor={Dev Goyal}}
\title{A396093\\[3pt]A complete parity classification}
\author{Dev Goyal}
\date{24 October 2026}
\begin{document}
\maketitle
\begin{center}
\fbox{\begin{minipage}{0.91\linewidth}\small
\textbf{AI contribution and review status.}
ChatGPT (OpenAI) generated the mathematical exposition, proof development,
Lean interfaces, and checks for this collection in response to Dev Goyal's
prompts. Dev Goyal is the sole human author. Independent human review is
pending. Original sequence authors, known mathematical results, and adapted
upstream code retain their separate credits. This document makes no claim
of publication, OEIS acceptance, or worldwide priority.
\end{minipage}}
\end{center}
\begin{abstract}
The third compositional iterate of $x/(1-x)^2$ has odd coefficients exactly at indices congruent to $1,3,7,9$ modulo ten. We prove this full classification and both parity assertions recorded in the source entry. The proof reduces an actual integral rational composition to a characteristic-two denominator identity.
\end{abstract}
\section{A396093: a complete parity classification}

The original definition is
\[
B(x)=\frac{x}{(1-x)^2},\qquad A=B\circ B\circ B,
\qquad a_n=[x^n]A\in\Z.
\]
The offset is zero, so $a_0=0$. These definitions, rather than a conjectured
coefficient pattern, are the starting point of the Lean construction.

\begin{theorem}\label{thm:parity}
For every $n\ge0$,
\[
a_n\equiv
\begin{cases}
1\pmod2,&n\equiv1,3,7,9\pmod{10},\\
0\pmod2,&\text{otherwise}.
\end{cases}
\]
Consequently $a_{2n}$ is even for every $n\ge0$. For $n\ge1$,
\[
a_{2n-1}\text{ is even}
\quad\Longleftrightarrow\quad n\equiv3\pmod5
\quad\Longleftrightarrow\quad n=5k-2\text{ for some }k\ge1.
\]
\end{theorem}

\begin{proof}
Reduce coefficients to $\F_2$. In characteristic two,
$B(x)=x/(1+x^2)$. If $T=P/D$ has zero constant coefficient and $D(0)=1$,
then
\[
B(T)=\frac{PD}{D^2+P^2}.
\]
Every denominator has constant coefficient one and is therefore an
invertible formal series. Applying this identity gives
\[
\begin{array}{c|c|c}
j & \text{numerator of }B^{\circ j} & \text{denominator}\\ \hline
1 & x & 1+x^2\\
2 & x+x^3 & 1+x^2+x^4\\
3 & x+x^7 & 1+x^2+x^4+x^6+x^8.
\end{array}
\]
Multiplication of the last numerator and denominator by $1+x^2$ gives
the exact identity
\begin{equation}\label{eq:parityrational}
A(x)(1+x^{10})=x+x^3+x^7+x^9\qquad\text{in }\ps{\F_2}.
\end{equation}
Since $(1+x^{10})^{-1}=\sum_{j\ge0}x^{10j}$ in this ring,
\eqref{eq:parityrational} yields precisely the four stated residue classes.
Substitute $2n$ and $2n-1$ to obtain the two original assertions. For
the latter, $2n-1$ is the omitted odd residue class $5$ modulo ten
exactly when $n\equiv3\pmod5$.
\end{proof}

The formal proof works with the integral series $x(\sum_{n\ge0}x^n)^2$
and its actual third composition. It proves that reduction commutes with
composition. The denominator cancellation in
\eqref{eq:parityrational} uses a constructed unit inverse; it does not
assume cancellation in an arbitrary ring. The final Lean statements
retain the source's $n\ge1$ and $k\ge1$ restrictions where subtraction
of natural numbers occurs.

\section{Lean proof and exact scope}
The dedicated entry file is \path{lean/Entries/A396093.lean}. Its declarations lie in namespace \texttt{FinalCollection.A396093}. The principal checked endpoints are:
\begin{itemize}[nosep]
\item \texttt{source\_denominator}: the literal rational denominator identity.
\item \texttt{source\_composition}: the threefold composition defining the sequence.
\item \texttt{reduced\_denominator}: the characteristic-two period identity.
\item \texttt{complete\_parity}: the all-index residue classification.
\item \texttt{even\_iff}: the equivalent integer-evenness characterization.
\item \texttt{even\_indices}: the first source assertion, including index zero.
\item \texttt{odd\_indices}: the exact positive-index source formulation.
\end{itemize}
The full implementations are in \path{lean/More/RationalParity.lean}.
The source offset is zero. The coefficient statement concerns the actual third iterate, not a separately postulated periodic sequence. The two source assertions were labeled conjectures in the retrieved entry; the full result here includes every $n\ge0$.

The companion is a separate, explicitly typed Lean entry file. Its proofs
invoke the preserved implementation modules; the complete dependency source
is included in the same archive. Shared lemmas are reused rather than
assumed. The entry file is not a self-contained replacement for Mathlib or
the supplied common library.

The project pins Lean 4.19.0 and Mathlib revision
\begin{center}\small\ttfamily
c44e0c8ee63ca166450922a373c7409c5d26b00b.
\end{center}
The collection checker recompiles the proof sources and audits the
theorems' transitive axiom dependencies. Only \texttt{propext},
\texttt{Classical.choice}, and \texttt{Quot.sound} are permitted.
The delivered verification reports record the actual successful runs.
Finite numerical checks are corroboration; the displayed conclusions are
unbounded theorems. An AI review of source correspondence does not replace
independent human review.
\begin{thebibliography}{9}
\raggedright
\bibitem{oeis} OEIS Foundation, entry A396093,
\url{https://oeis.org/A396093}. Source checked 24 October 2026.
The source sequence and selected conjectures are credited to Paul D. Hanna.
\bibitem{repo} D. Goyal, supplied reference repository,
\url{https://github.com/dvgyl/oeis-a396843-proof}.
This separate presentation belongs to the final collected proof package;
no earlier DOI is reassigned to it.
\end{thebibliography}
\end{document}
